\section{No-go theorem of internal interpretation}
\label{no-go-section}

So we want to build a reflective type system. That is,
a system where a user is able to both define a model of
a system internally and apply it to well-typed terms to
get their internal interpretations. Naturally, we would
like user-defined models to respect definitional
equalities of a system; in a dependently typed system,
this could be naturally achieved by using an identity
type to spell out corresponding equalities inside a
model. Then, it would be nice if, before interpreting
a term, system would first bring it to normal form to
get rid of the artifacts of the construction in the
result of the interpretation. Finally, in a dependently
typed system, it is natural to expect reductions of open
terms, so we would like our interpretation to compute on
open terms.

The theorem proven below asserts that we cannot have all
three in a single system.

First, a couple of definitions. We're going to work
within a classical type-theoretical framework, which
introduces a sort of contexts $\Gamma$, terms $t,t'$ and
types $T$ linked by the typing relation
$\Gamma\vdash t:T$; computation is specified by a
one-step reduction relation $t\leadsto t'$ on terms.

\begin{definition}
    [Na\"ive computational interpretation]
    We call a type operator $F[\_]$ a na\"ive
    computational interpretation of types if it
    respects computation on types (that is, if there is
    a reduction relation on types $(\leadsto)$ and
    $T\leadsto T'$, then $F[T]\leadsto^* F[T']$), and
    computes as a primitive recursive function on normal
    forms of types. Analogously, term with a hole
    $\llbracket\_\rrbracket$ is a na\"ive
    computational interpretation of terms if it
    respects computation on terms, computes as a
    primitive recursive function on normal forms of
    terms and, if $\Gamma\vdash t:T$, then
    $\Gamma\vdash\llbracket t\rrbracket:F[T]$.
\end{definition}

\begin{theorem}[No-go theorem of naive interpretation]
\label{no-go}
    In a confluent type system with a type of functions
    and na\"{\i}ve computational interpretation, if
    reduction commutes with substitution, for arbitrary
    well-typed $s$ and $t$ there exists $u$ such that:
    \[
        \mathtt{app}\llbracket\lambda x.t\rrbracket
        \llbracket s \rrbracket \leadsto^* u\;
        \text{\reflectbox{$\leadsto^*$}}\;
        \llbracket t[s/x] \rrbracket
    \]
    where $\mathtt{app}$ is how $\llbracket\_\rrbracket$
    acts on applications.
\end{theorem}

Therefore, $\mathtt{app}\llbracket\lambda x.t\rrbracket
\llbracket s \rrbracket$ and $\llbracket t[s/x]
\rrbracket$ are definitionally equivalent, which forces
beta-reductions in a model to hold definitionally rather
than propositionally, directly contradicting our desired
outcome. Analogous theorems can be proven for any other
reduction present in the system; in other words,
\textbf{propositional equality is not strong enough}.
The proof is by direct example:

\begin{proof}
    We are given well-typed $\Gamma\vdash s:S$ and
    $\Gamma,x:S\vdash t:T$. Consider the following term:
    \[
        \Gamma,x:S\to T,y:S\vdash
            \llbracket x\,y\rrbracket:F[T]
    \]
    Term under interpretation is in normal form, hence
    the whole term computes to $\mathtt{app}\llbracket x
    \rrbracket\llbracket y\rrbracket$.
    Substituting $\lambda x.t$ for $x$ and $s$ for $y$,
    we get $\mathtt{app}\llbracket\lambda x.t\rrbracket
    \llbracket s \rrbracket$.

    On the other hand, if we perform substitution first,
    we would get $\Gamma\vdash\llbracket(\lambda x.t)s
    \rrbracket:F[T]$. Since na\"{\i}ve interpretation
    respects reduction, this reduces to
    $\llbracket t[s/x] \rrbracket$.

    Since reduction commutes with substitution and
    system is confluent, there should be such $u$ that
    both terms reduce to.
\end{proof}

Fundamentally, the proof above exploits the openness of
the original term; if we restrict interpretation to
compute on closed terms only, we can arrive at a sound
system without requiring definitional equalities in
models, as exemplified by our non-dependently typed
system $\SSTLC$ described in Section~\ref{SSTLC}. And
this, actually, is a well-known setup: under the same
conditions (no dependent types + computation on closed
terms), self-interpreter for System~F$\omega$ was
developed \cite{brown2016breaking}.

For a dependently typed system, however, our no-go
theorem leaves us with three options:

\begin{enumerate}
    \item Restrict models to those where equalities hold
        definitionally rather than propositionally.
    \item Do not proceed with interpretation unless
        the term is closed. To allow reasoning on open
        terms, equations describing the behavior of
        interpretation are exposed as propositional
        axioms.
    \item Do not reduce open terms under
        interpretation at all, performing the
        interpretation as-is. To simplify the result
        of interpretation, one can rewrite it with
        equalities present in the model.
\end{enumerate}

If you view existing work on displacement
\cite{McBride:Crude2011, McBride:Crude2012}, internal
parametricity \cite{Altenkirch:Parametricity2024} and
displayed type theory \cite{Kolomatskaia_Shulman_2025}
through the lens of interpretation inside a fixed
model (or family of models), you'll see that
corresponding models indeed validate equalities
definitionally; thus it seems that the first approach
needs to be taken. It is also the most satisfying from
the computational point of view. However, if you want a
sound system where models can be user-defined, you need
to constrain models through typing discipline. By
following techniques currently known to us
\cite{bezem2021notegeneralizedalgebraictheories},
constraining models in such a way would require
defining a type of models as a type of algebras of an
inductive-inductive-recursive type family (or even
inductive-inductive-recursive simplicial types, for
univalent type theories) which is well beyond the scope
of this paper.

Out of the two remaining options, the last one is better
suited for practical dependently-typed systems since it
properly computes on open terms. In addition, it doesn't
restrict the class of possible models while at the same
time providing the benefits of the first approach: if
a model is such that the equations hold definitionally,
no rewriting on the result of interpretation is needed
to reduce it. The only downside is that, possibly,
reductions \textit{after} interpretation are more costly
than the ones we could perform \textit{before} it.
